\documentclass[12pt, a4paper]{article}

\usepackage{amsmath, amssymb, amsthm}
\usepackage{geometry}
\usepackage{hyperref}
\usepackage{algorithm}
\usepackage{algpseudocode}
\usepackage{booktabs}

\geometry{margin=2.5cm}

\title{%
  \textbf{STAS Simulation Mathematics}\\[0.5em]
  \large Closeness Centrality, Monte Carlo Estimation,\\
  and Confidence Interval Construction
}
\author{Socio-Technical Alignment Simulator}
\date{Sprint 4 --- May 2026}

\newtheorem{definition}{Definition}
\newtheorem{proposition}{Proposition}
\newtheorem{remark}{Remark}

% ──────────────────────────────────────────────────────────────────────────────
\begin{document}
\maketitle
\tableofcontents
\newpage

% ─────────────────────────────────────────────────────────────────────────────
\section{Graph Model}

Let $G = (V, E, w)$ be a weighted, undirected knowledge graph where:

\begin{itemize}
  \item $V$ is the set of \textbf{Engineer} nodes with $|V| = n$.
  \item $E \subseteq V \times V$ is the edge set representing observed
        collaboration events (PR reviews, Slack threads, Jira blockers,
        pair-programming sessions).
  \item $w : E \to \mathbb{Z}_{>0}$ is a weight function representing the
        count of collaboration events between two engineers over a trailing
        90-day window.
\end{itemize}

Candidate insertion extends $G$ to $G' = (V \cup \{v_c\}, E', w')$, where
$v_c$ is the candidate node and $E'$ contains the original edges plus
projected edges
\[
  (v_c,\, u) \in E' \iff |\text{skills}(v_c) \cap \text{skills}(u)| > 0,
\]
with weight $w'(v_c, u) = |\text{skills}(v_c) \cap \text{skills}(u)|$.

% ─────────────────────────────────────────────────────────────────────────────
\section{Closeness Centrality}

\begin{definition}[Geodesic distance]
For nodes $u, v \in V$, the geodesic distance $d(u,v)$ is the length of the
shortest path between them in $G$, where path length is the number of edges
traversed (unweighted BFS).  If no path exists, $d(u,v) = \infty$.
\end{definition}

\begin{definition}[Closeness centrality --- Wasserman-Faust normalisation]
Let $R(v) = \{u \in V : u \neq v,\, d(u,v) < \infty\}$ be the set of nodes
reachable from $v$, and let $n_r = |R(v)|$.  The closeness centrality of $v$
in $G$ is

\begin{equation}
  \label{eq:cc}
  C_C(v) =
  \frac{n_r}{\displaystyle\sum_{u \in R(v)} d(u,v)}
  \cdot
  \frac{n_r - 1}{n - 1}
\end{equation}

The second factor is the Wasserman-Faust correction ensuring $C_C(v) \leq 1$
for all graphs, including disconnected ones.
\end{definition}

\begin{proposition}[Range]
$0 \leq C_C(v) \leq 1$ for all $v \in V$.
\end{proposition}
\begin{proof}
The first factor $n_r / \sum d(u,v)$ satisfies $0 \leq \cdot \leq 1$ because
$\sum d(u,v) \geq n_r$ (each reachable node is at distance $\geq 1$).  The
second factor $(n_r - 1)/(n-1) \leq 1$ by definition of $R(v)$.
For an isolated node $R(v) = \emptyset$, we define $C_C(v) = 0$.
\end{proof}

\begin{proposition}[Universal hub]
If $v$ is adjacent to all other nodes and all shortest paths have length 1,
then $C_C(v) = 1$.
\end{proposition}
\begin{proof}
$n_r = n-1$, $\sum d(u,v) = n-1$, so first factor $= 1$.
Second factor $(n-1-1)/(n-1) \neq 1$ in general --- the hub achieves $C_C = 1$
only when all other nodes are \emph{directly connected} to it (distance 1) and
the graph is connected.  In this case first factor $= (n-1)/(n-1) = 1$ and
second factor $= (n-1)/(n-1) = 1$.
\end{proof}

\begin{remark}[Implementation]
Equation~\ref{eq:cc} is implemented in \texttt{simulation/metrics.py::closeness\_centrality}
using \texttt{networkx.single\_source\_shortest\_path\_length} for the BFS from $v$.
Graphs with $|V| \leq 200$ use NetworkX; larger graphs use the Neo4j GDS
\texttt{gds.closeness.stream()} procedure.
\end{remark}

% ─────────────────────────────────────────────────────────────────────────────
\section{Additional Metrics}

\subsection{Betweenness Centrality Delta}

The betweenness centrality of node $u$ is
\[
  B(u) = \sum_{s \neq u \neq t} \frac{\sigma_{st}(u)}{\sigma_{st}}
\]
where $\sigma_{st}$ is the number of shortest paths from $s$ to $t$ and
$\sigma_{st}(u)$ is the number of those paths passing through $u$.

The \emph{betweenness centrality delta} measures how much candidate insertion
changes the mean betweenness of existing engineers:
\[
  \Delta B = \bar{B}(G'_{-v_c}) - \bar{B}(G \setminus \{v_c\})
\]
where $\bar{B}$ denotes the mean over all non-candidate nodes.
A negative $\Delta B$ indicates that $v_c$ reduces existing bottlenecks.

\subsection{Silo Risk Score}

The silo risk score uses the clustering coefficient:
\[
  C(u) = \frac{|\{(v,w) \in E : v, w \in N(u)\}|}{k_u (k_u - 1) / 2}
\]
where $N(u)$ is the neighbourhood of $u$ and $k_u = |N(u)|$.

\[
  \text{SiloRisk}(v_c) = \bar{C}(G') - \bar{C}(G)
\]
\begin{itemize}
  \item Negative: candidate bridges clusters, reducing silos.
  \item Positive: candidate reinforces a single cluster, increasing silos.
\end{itemize}

\subsection{Time-to-Value Estimate}

A calibrated heuristic mapping closeness centrality to integration time:
\[
  \text{TTV}(v_c) = \frac{2.0}{\max\!\left(C_C(v_c),\; \tfrac{1}{26}\right)} \;\; \text{weeks}
\]
The $1/26$ floor caps the estimate at 52 weeks (one year) for near-isolated
candidates.

\begin{center}
\begin{tabular}{ccc}
\toprule
$C_C(v_c)$ & TTV (weeks) & Interpretation \\
\midrule
1.00 & 2  & Immediate knowledge-hub integration \\
0.50 & 4  & Moderate onboarding path \\
0.25 & 8  & Extended integration period \\
0.10 & 20 & High isolation risk \\
0.00 & 52 & Integration unlikely within one year \\
\bottomrule
\end{tabular}
\end{center}

% ─────────────────────────────────────────────────────────────────────────────
\section{Monte Carlo Simulation}

\subsection{Stochastic Edge Weight Model}

Real collaboration counts fluctuate across 90-day windows.  We model this
uncertainty by sampling each edge weight from a Poisson distribution:
\[
  w^{(k)}(e) \sim \text{Poisson}\!\left(\lambda = w(e)\right), \quad
  w^{(k)}(e) \geq 1
\]
where $w(e)$ is the observed base weight.  The Poisson distribution is the
natural choice because edge weights count discrete collaboration events.

\subsection{Algorithm}

\begin{algorithm}
\caption{STAS Monte Carlo Simulation}
\begin{algorithmic}[1]
\Require Snapshot $G=(V,E,w)$, candidate $v_c$, $N$ iterations, seed $s$
\Ensure SimulationResult with distributions for all metrics
\State $\text{rng} \leftarrow \text{Generator}(s)$ \Comment{Seeded RNG for determinism}
\State Build $G_0$ from snapshot; add $v_c$ with projected edges
\State $\mathbf{cc}, \mathbf{sr} \leftarrow$ empty arrays of size $N$
\For{$k = 1$ to $N$}
  \State $G^{(k)} \leftarrow \text{copy}(G_0)$
  \For{each edge $e \in G^{(k)}$}
    \State $w^{(k)}(e) \leftarrow \max(1, \text{rng.Poisson}(w_0(e)))$
  \EndFor
  \State $\mathbf{cc}[k] \leftarrow C_C(G^{(k)}, v_c)$ \quad \emph{(BFS from $v_c$)}
  \State $\mathbf{sr}[k] \leftarrow \bar{C}(G^{(k)}) - \bar{C}(G^{(k)} \setminus \{v_c\})$
\EndFor
\State $\Delta B \leftarrow$ betweenness delta on mean-weight graph $\bar{G}$
\State \Return aggregate statistics for all sampled quantities
\end{algorithmic}
\end{algorithm}

\subsection{Confidence Interval Construction}

For a sample $\{x_1, \ldots, x_N\}$, the 95\% confidence interval uses the
normal approximation (valid for $N \geq 30$):
\[
  \text{CI}_{95} = \bar{x} \pm 1.96 \cdot \frac{s}{\sqrt{N}}
\]
where $\bar{x}$ is the sample mean and $s$ is the unbiased sample standard
deviation ($\text{ddof}=1$).

\subsection{Determinism Guarantee}

The simulation is \emph{byte-identical} for the same inputs $(G, v_c, N, s)$
because:
\begin{enumerate}
  \item The RNG is \texttt{numpy.random.default\_rng(seed)} (PCG64).
  \item No external I/O, OS entropy, or wall-clock time is used during the run.
  \item The graph copy operation is deterministic for the same adjacency list.
\end{enumerate}

% ─────────────────────────────────────────────────────────────────────────────
\section{Performance Analysis}

\begin{center}
\begin{tabular}{lll}
\toprule
Operation & Complexity & Notes \\
\midrule
Build base graph & $O(n + m)$ & Once per run \\
Graph copy (per iter) & $O(n + m)$ & NetworkX \texttt{copy()} \\
Closeness BFS (per iter) & $O(n + m)$ & Single source \\
Clustering coefficient (per iter) & $O(n \bar{k}^2)$ & $\bar{k}$ = mean degree \\
Betweenness delta & $O(nm)$ & Once (post-loop) \\
Total for $N$ iterations & $O(N(n+m))$ & Dominated by loop \\
\bottomrule
\end{tabular}
\end{center}

For $n=50$, $m=490$ ($40\%$ density), $N=1{,}000$:
total $\approx O(540{,}000)$ graph traversal operations.
Empirical baseline: $< 2$ seconds in CPython 3.11 on commodity hardware.

\end{document}
